% !TeX encoding = UTF-8
% !TeX program = xelatex
% !TeX spellcheck = en_US

\documentclass[degree=master]{thuthesis}
  % 学位 degree:
  %   doctor | master | bachelor | postdoc
  % 学位类型 degree-type:
  %   academic（默认）| professional
  % 语言 language
  %   chinese（默认）| english
  % 字体库 fontset
  %   windows | mac | fandol | ubuntu
  % 建议终版使用 Windows 平台的字体编译


% 论文基本配置，加载宏包等全局配置
% !TeX root = ./thuthesis-example.tex

% 论文基本信息配置

\thusetup{
  %******************************
  % 注意：
  %   1. 配置里面不要出现空行
  %   2. 不需要的配置信息可以删除
  %   3. 建议先阅读文档中所有关于选项的说明
  %******************************
  %
  % 输出格式
  %   选择打印版（print）或用于提交的电子版（electronic），前者会插入空白页以便直接双面打印
  %
  output = print,
  % 格式类型
  %   默认为论文（thesis），也可以设置为开题报告（proposal）
  % thesis-type = proposal,
  %
  % 标题
  %   可使用“\\”命令手动控制换行
  %
  title  = {清华大学学位论文 \LaTeX{} 模板\\使用示例文档 v\version},
  title* = {An Introduction to \LaTeX{} Thesis Template of Tsinghua
            University v\version},
  %
  % 学科门类
  %   1. 学术型
  %      - 中文
  %        需注明所属的学科门类，例如：
  %        哲学、经济学、法学、教育学、文学、历史学、理学、工学、农学、医学、
  %        军事学、管理学、艺术学
  %      - 英文
  %        博士：Doctor of Philosophy
  %        硕士：
  %          哲学、文学、历史学、法学、教育学、艺术学门类，公共管理学科
  %          填写“Master of Arts“，其它填写“Master of Science”
  %   2. 专业型
  %      直接填写专业学位的名称，例如：
  %      教育博士、工程硕士等
  %      Doctor of Education, Master of Engineering
  %   3. 本科生不需要填写
  %
  degree-category  = {工学硕士},
  degree-category* = {Master of Science},
  %
  % 培养单位
  %   填写所属院系的全名
  %
  department = {计算机科学与技术系},
  %
  % 学科
  %   1. 研究生学术型学位，获得一级学科授权的学科填写一级学科名称，其他填写二级学科名称
  %   2. 本科生填写专业名称，第二学位论文需标注“（第二学位）”
  %
  discipline  = {计算机科学与技术},
  discipline* = {Computer Science and Technology},
  %
  % 专业领域
  %   1. 设置专业领域的专业学位类别，填写相应专业领域名称
  %   2. 2019 级及之前工程硕士学位论文，在 `engineering-field` 填写相应工程领域名称
  %   3. 其他专业学位类别的学位论文无需此信息
  %
  % professional-field  = {计算机技术},
  % professional-field* = {Computer Technology},
  %
  % 姓名
  %
  author  = {薛瑞尼},
  author* = {Xue Ruini},
  %
  % 学号
  % 仅当书写开题报告时需要（同时设置 `thesis-type = proposal'）
  %
  % student-id = {2000310000},
  %
  % 指导教师
  %   中文姓名和职称之间以英文逗号“,”分开，下同
  %
  supervisor  = {郑纬民, 教授},
  supervisor* = {Professor Zheng Weimin},
  %
  % 副指导教师
  %
  associate-supervisor  = {陈文光, 教授},
  associate-supervisor* = {Professor Chen Wenguang},
  %
  % 联合指导教师
  %
  % co-supervisor  = {某某某, 教授},
  % co-supervisor* = {Professor Mou Moumou},
  %
  % 日期
  %   使用 ISO 格式；默认为当前时间
  %
  % date = {2019-07-07},
  %
  % 是否在中文封面后的空白页生成书脊（默认 false）
  %
  include-spine = false,
  %
  % 密级和年限
  %   秘密, 机密, 绝密
  %
  % secret-level = {秘密},
  % secret-year  = {10},
  %
  % 博士后专有部分
  %
  % clc                = {分类号},
  % udc                = {UDC},
  % id                 = {编号},
  % discipline-level-1 = {计算机科学与技术},  % 流动站（一级学科）名称
  % discipline-level-2 = {系统结构},          % 专业（二级学科）名称
  % start-date         = {2011-07-01},        % 研究工作起始时间
}

% 载入所需的宏包

% 定理类环境宏包
\usepackage{amsthm}
% 也可以使用 ntheorem
% \usepackage[amsmath,thmmarks,hyperref]{ntheorem}

\thusetup{
  %
  % 数学字体
  % math-style = GB,  % GB | ISO | TeX
  math-font  = xits,  % stix | xits | libertinus
}

% 可以使用 nomencl 生成符号和缩略语说明
% \usepackage{nomencl}
% \makenomenclature

% 表格加脚注
\usepackage{threeparttable}

% 表格中支持跨行
\usepackage{multirow}

% 固定宽度的表格。
% \usepackage{tabularx}

% 跨页表格
\usepackage{longtable}

% 算法
\usepackage{algorithm}
\usepackage{algorithmic}

% 量和单位
\usepackage{siunitx}

% 参考文献使用 BibTeX + natbib 宏包
% 顺序编码制
\usepackage[sort]{natbib}
\bibliographystyle{thuthesis-numeric}

% 著者-出版年制
% \usepackage{natbib}
% \bibliographystyle{thuthesis-author-year}

% 生命科学学院要求使用 Cell 参考文献格式（2023 年以前使用 author-date 格式）
% \usepackage{natbib}
% \bibliographystyle{cell}

% 本科生参考文献的著录格式
% \usepackage[sort]{natbib}
% \bibliographystyle{thuthesis-bachelor}

% 参考文献使用 BibLaTeX 宏包
% \usepackage[style=thuthesis-numeric]{biblatex}
% \usepackage[style=thuthesis-author-year]{biblatex}
% \usepackage[style=gb7714-2015]{biblatex}
% \usepackage[style=apa]{biblatex}
% \usepackage[style=mla-new]{biblatex}
% 声明 BibLaTeX 的数据库
% \addbibresource{ref/refs.bib}

% 定义所有的图片文件在 figures 子目录下
\graphicspath{{figures/}}

% 数学命令
\makeatletter
\newcommand\dif{%  % 微分符号
  \mathop{}\!%
  \ifthu@math@style@TeX
    d%
  \else
    \mathrm{d}%
  \fi
}
\makeatother

% hyperref 宏包在最后调用
\usepackage{hyperref}



\begin{document}

% 封面
\maketitle

% 学位论文指导小组、公开评阅人和答辩委员会名单
% 本科生不需要
% !TeX root = ../thuthesis-example.tex

\begin{committee}[name={学位论文指导小组、公开评阅人和答辩委员会名单}]

  \newcolumntype{C}[1]{@{}>{\centering\arraybackslash}p{#1}}

  \section*{指导小组名单}

  \begin{center}
    \begin{tabular}{C{3cm}C{3cm}C{9cm}@{}}
      李XX & 教授     & 清华大学 \\
      王XX & 副教授   & 清华大学 \\
      张XX & 助理教授 & 清华大学 \\
    \end{tabular}
  \end{center}


  \section*{公开评阅人名单}

  \begin{center}
    \begin{tabular}{C{3cm}C{3cm}C{9cm}@{}}
      刘XX & 教授   & 清华大学                    \\
      陈XX & 副教授 & XXXX大学                    \\
      杨XX & 研究员 & 中国XXXX科学院XXXXXXX研究所 \\
    \end{tabular}
  \end{center}


  \section*{答辩委员会名单}

  \begin{center}
    \begin{tabular}{C{2.75cm}C{2.98cm}C{4.63cm}C{4.63cm}@{}}
      主席 & 赵XX                  & 教授                    & 清华大学       \\
      委员 & 刘XX                  & 教授                    & 清华大学       \\
          & \multirow{2}{*}{杨XX} & \multirow{2}{*}{研究员} & 中国XXXX科学院 \\
          &                       &                         & XXXXXXX研究所  \\
          & 黄XX                  & 教授                    & XXXX大学       \\
          & 周XX                  & 副教授                  & XXXX大学       \\
      秘书 & 吴XX                  & 助理研究员              & 清华大学       \\
    \end{tabular}
  \end{center}

\end{committee}



% 也可以导入 Word 版转的 PDF 文件
% \begin{committee}[file=figures/committee.pdf]
% \end{committee}


% 使用授权的说明
% 本科生开题报告不需要
\copyrightpage
% 将签字扫描后授权文件 scan-copyright.pdf 替换原始页面
% \copyrightpage[file=scan-copyright.pdf]

\frontmatter
\setlength{\baselineskip}{20pt}
\begin{abstract}
	[内容为小四号宋体。] 中文摘要应将学位论文的内容要点简短明了地表达出来，硕士学位论文一般为500~1000字，博士学位论文一般为1000~2000字。留学生英文版学位论文不少于3000字中文摘要，留学生英文版博士学位论文不少于5000字中文摘要。字体为宋体小四号。内容应包括工作目的、研究方法、成果和结论。要突出本论文的创新点，语言力求精炼。为了便于文献检索，应在本页下方另起一行注明论文的关键词（3-8个），如有可能，尽量采用《汉语主题词表》等词表提供的规范词。图X幅，表X个，参考文献X篇。

		
	\noindent\keywords{[请输入关键词（3-8），以分号分隔。] }
\end{abstract}

% 目录
\tableofcontents

% 插图和附表清单
% 本科生的插图索引和表格索引需要移至正文之后、参考文献前
% \listoffiguresandtables  % 插图和附表清单（仅限研究生）
\listoffigures           % 插图清单
\listoftables            % 附表清单

% 符号对照表
% !TeX root = ../thuthesis-example.tex

\begin{denotation}[3cm]
  \item[PI] 聚酰亚胺
  \item[MPI] 聚酰亚胺模型化合物，N-苯基邻苯酰亚胺
  \item[PBI] 聚苯并咪唑
  \item[MPBI] 聚苯并咪唑模型化合物，N-苯基苯并咪唑
  \item[PY] 聚吡咙
  \item[PMDA-BDA] 均苯四酸二酐与联苯四胺合成的聚吡咙薄膜
  \item[MPY] 聚吡咙模型化合物
  \item[As-PPT] 聚苯基不对称三嗪
  \item[MAsPPT] 聚苯基不对称三嗪单模型化合物，3,5,6-三苯基-1,2,4-三嗪
  \item[DMAsPPT] 聚苯基不对称三嗪双模型化合物（水解实验模型化合物）
  \item[S-PPT] 聚苯基对称三嗪
  \item[MSPPT] 聚苯基对称三嗪模型化合物，2,4,6-三苯基-1,3,5-三嗪
  \item[PPQ] 聚苯基喹噁啉
  \item[MPPQ] 聚苯基喹噁啉模型化合物，3,4-二苯基苯并二嗪
  \item[HMPI] 聚酰亚胺模型化合物的质子化产物
  \item[HMPY] 聚吡咙模型化合物的质子化产物
  \item[HMPBI] 聚苯并咪唑模型化合物的质子化产物
  \item[HMAsPPT] 聚苯基不对称三嗪模型化合物的质子化产物
  \item[HMSPPT] 聚苯基对称三嗪模型化合物的质子化产物
  \item[HMPPQ] 聚苯基喹噁啉模型化合物的质子化产物
  \item[PDT] 热分解温度
  \item[HPLC] 高效液相色谱（High Performance Liquid Chromatography）
  \item[HPCE] 高效毛细管电泳色谱（High Performance Capillary lectrophoresis）
  \item[LC-MS] 液相色谱-质谱联用（Liquid chromatography-Mass Spectrum）
  \item[TIC] 总离子浓度（Total Ion Content）
  \item[\textit{ab initio}] 基于第一原理的量子化学计算方法，常称从头算法
  \item[DFT] 密度泛函理论（Density Functional Theory）
  \item[$E_a$] 化学反应的活化能（Activation Energy）
  \item[ZPE] 零点振动能（Zero Vibration Energy）
  \item[PES] 势能面（Potential Energy Surface）
  \item[TS] 过渡态（Transition State）
  \item[TST] 过渡态理论（Transition State Theory）
  \item[$\increment G^\neq$] 活化自由能（Activation Free Energy）
  \item[$\kappa$] 传输系数（Transmission Coefficient）
  \item[IRC] 内禀反应坐标（Intrinsic Reaction Coordinates）
  \item[$\nu_i$] 虚频（Imaginary Frequency）
  \item[ONIOM] 分层算法（Our own N-layered Integrated molecular Orbital and molecular Mechanics）
  \item[SCF] 自洽场（Self-Consistent Field）
  \item[SCRF] 自洽反应场（Self-Consistent Reaction Field）
\end{denotation}



% 也可以使用 nomencl 宏包，需要在导言区
% \usepackage{nomencl}
% \makenomenclature

% 在这里输出符号说明
% \printnomenclature[3cm]

% 在正文中的任意为都可以标题
% \nomenclature{PI}{聚酰亚胺}
% \nomenclature{MPI}{聚酰亚胺模型化合物，N-苯基邻苯酰亚胺}
% \nomenclature{PBI}{聚苯并咪唑}
% \nomenclature{MPBI}{聚苯并咪唑模型化合物，N-苯基苯并咪唑}
% \nomenclature{PY}{聚吡咙}
% \nomenclature{PMDA-BDA}{均苯四酸二酐与联苯四胺合成的聚吡咙薄膜}
% \nomenclature{MPY}{聚吡咙模型化合物}
% \nomenclature{As-PPT}{聚苯基不对称三嗪}
% \nomenclature{MAsPPT}{聚苯基不对称三嗪单模型化合物，3,5,6-三苯基-1,2,4-三嗪}
% \nomenclature{DMAsPPT}{聚苯基不对称三嗪双模型化合物（水解实验模型化合物）}
% \nomenclature{S-PPT}{聚苯基对称三嗪}
% \nomenclature{MSPPT}{聚苯基对称三嗪模型化合物，2,4,6-三苯基-1,3,5-三嗪}
% \nomenclature{PPQ}{聚苯基喹噁啉}
% \nomenclature{MPPQ}{聚苯基喹噁啉模型化合物，3,4-二苯基苯并二嗪}
% \nomenclature{HMPI}{聚酰亚胺模型化合物的质子化产物}
% \nomenclature{HMPY}{聚吡咙模型化合物的质子化产物}
% \nomenclature{HMPBI}{聚苯并咪唑模型化合物的质子化产物}
% \nomenclature{HMAsPPT}{聚苯基不对称三嗪模型化合物的质子化产物}
% \nomenclature{HMSPPT}{聚苯基对称三嗪模型化合物的质子化产物}
% \nomenclature{HMPPQ}{聚苯基喹噁啉模型化合物的质子化产物}
% \nomenclature{PDT}{热分解温度}
% \nomenclature{HPLC}{高效液相色谱（High Performance Liquid Chromatography）}
% \nomenclature{HPCE}{高效毛细管电泳色谱（High Performance Capillary lectrophoresis）}
% \nomenclature{LC-MS}{液相色谱-质谱联用（Liquid chromatography-Mass Spectrum）}
% \nomenclature{TIC}{总离子浓度（Total Ion Content）}
% \nomenclature{\textit{ab initio}}{基于第一原理的量子化学计算方法，常称从头算法}
% \nomenclature{DFT}{密度泛函理论（Density Functional Theory）}
% \nomenclature{$E_a$}{化学反应的活化能（Activation Energy）}
% \nomenclature{ZPE}{零点振动能（Zero Vibration Energy）}
% \nomenclature{PES}{势能面（Potential Energy Surface）}
% \nomenclature{TS}{过渡态（Transition State）}
% \nomenclature{TST}{过渡态理论（Transition State Theory）}
% \nomenclature{$\increment G^\neq$}{活化自由能（Activation Free Energy）}
% \nomenclature{$\kappa$}{传输系数（Transmission Coefficient）}
% \nomenclature{IRC}{内禀反应坐标（Intrinsic Reaction Coordinates）}
% \nomenclature{$\nu_i$}{虚频（Imaginary Frequency）}
% \nomenclature{ONIOM}{分层算法（Our own N-layered Integrated molecular Orbital and molecular Mechanics）}
% \nomenclature{SCF}{自洽场（Self-Consistent Field）}
% \nomenclature{SCRF}{自洽反应场（Self-Consistent Reaction Field）}



% 正文部分
\mainmatter

%%% Local Variables:
%%% mode: latex
%%% TeX-master: t
%%% End:

\chapter{����}
\label{cha:intro}
[����������ѡ��ö��䣬�����滻֮������ΪС�ĺ����塣] ���ԣ������ۣ���Ҫ˵���о�������Ŀ�ġ���Χ����������ǰ�˹�����֪ʶ�հס����ۻ����ͷ������о����롢�о�������ʵ����ơ�Ԥ�ڽ��������ȡ�Ӧ�Լ����࣬��Ҫ��ժҪ��ͬ����Ҫ��ΪժҪ��ע�͡�һ��̿������е�֪ʶ���������в���׸����

\chapter{1������ �����룺\textbackslash chapter\{1������\}��}
[����������ѡ��ö��䣬�����滻֮������ΪС�ĺ����塣] ѧλ����Ϊ����Ҫ��ӳ������ȷ�������˼�ʵ�Ļ������ۺ�ϵͳ��ר��֪ʶ�����п����Ŀ�ѧ��Ұ�����о��������˳����֤����ˣ��й���ʷ�ع˺�ǰ�˹������ۺ��������Լ����۷����ȣ����Ե������£����㹻������������������ѧλ���ĵĺ��IJ��֣�ռ��Ҫƪ�������԰������������ʵ��͹۲ⷽ���������豸������ԭ�ϡ�ʵ��͹۲��������㷽���ͱ��ԭ�����������ϡ������ӹ�������ͼ�����γɵ��۵�͵����Ľ��۵ȡ�

�����о������漰��ѧ�ơ�ѡ�⡢�о��������������̡�������﷽ʽ���кܴ�IJ��죬���������ݲ�����ͳһ�Ĺ涨�����ǣ�����ʵ�����ǣ��͹����У�׼ȷ�걸���Ϻ��߼�����η����������ɶ���
\textbf{ͼ��}����������ͨ��ѧѧλ����׫д�淶��3.10.4

ͼӦ�б�š�ͼ�ı���ɡ�ͼ���ʹӡ�1����ʼ�İ�����������ɣ�ͼ�϶�ʱ���ɷ��������š�

ͼ����ͼ�⣬ͼ�⼴ͼ�����ƣ�����ͼ�ı��֮��ͼ�ı�ź�ͼ��Ӧ����ͼ�·���ͼ�������Ӣ�Ķ��գ�Ӣ�ģ�Times New Roman��������ţ�����������š�������д���������ϡ�

��ƬͼҪ���������Ҫ��ʾ���ֵ����������������ư档���÷Ŵ���С�ĸ���Ʒ�������������������С���Ƭ��Ӧ�б�ʾĿ����ߴ�ı�ȡ�

\textbf{����}����������ͨ��ѧѧλ����׫д�淶��3.10.5

��Ӧ�б�š����ı���ɡ������ʹӡ�1����ʼ�İ�����������ɣ����϶�ʱ���ɷ��������š�

�����б��⣬���⼴�������ƣ����ڱ��ı��֮�󡣱��ı�źͱ���Ӧ���ڱ��Ϸ������������Ӣ�Ķ��գ����У��������ϡ�Ӣ�ģ�Times New Roman��������ţ�����������š�

���ı��ţ�һ�������ݺͲ�����Ŀ�������Һ������������������

���ı��Ž�����ù���ͨ�е����߱���

��ij������Ҫתҳ���ţ������ĸ�ҳ��Ӧ�ظ����ı�š���ź�����⣨��ʡ�ԣ��͡��������������ڱ��Ϸ���

������Ӧ�ظ���ͷ��

\textbf{��ʽ��}����������ͨ��ѧѧλ����׫д�淶��3.10.12

�����еĹ�ʽӦ�����𣬲�������д������Χ�������㹻�Ŀռ����ֿ���

�����������ϵĹ�ʽ��Ӧ�ôӡ�1����ʼ�İ��������ֽ��б�ţ�����������������ڡ���ʽ�ı���Ҷ˶��룬��ʽ����֮����á��������ӡ���ʽ�϶�ʱ���ɷ��±�š�

ʵ����

\begin{equation}
	{w_1} = {u_{11}} - {u_{12}}{u_{21}}
\end{equation}

Դ����Ϊ��
\begin{itemize}
	\item \textbackslash begin\{equation\}
	\item {w\_1} = {u\_{11}} - {u\_{12}}{u\_{21}}
	\item \textbackslash end\{equation\}
\end{itemize}

�ϳ��Ĺ�ʽ��Ҫת��ʱ��Ӧ�������ڡ����������У������ڡ�+��������������������/���ȼǺŴ����С���ʽ�з����ߵĺ��ߣ��䳤��Ӧ���ڻ��Դ��ڷ��Ӻͷ�ĸ�нϳ���һ����

����������д������Ӧ��������߶Ƚ���Ϊһ�С��罫��������дΪ��/���������Ÿ�Ϊ��ָ����

ʾ����

   ��$\frac{1}{{\sqrt 2 }}$�1/$\sqrt 2 $��${2^{ - 1/2}}$
   
\textbf{���ã�}����������ͨ��ѧѧλ����׫д�淶��3.10.7
���������õ����׵ı�ע��������GB/T 7714��2005������˳������ƣ����Բο�������ʽͳһ��š��������������׵ı�ʾӦ���������������һ���ֵ����Ͻǣ��������ױ���ð������������ڷ����š�[ ]���У���С4��������ϽDZꡣ

\textbf{ע�ͣ�}����������ͨ��ѧѧλ����׫д�淶��3.10.6
�������е��֡��ʻ�����Ҫ��һ������˵��������û�о����������Դʱ����ע�͡�Ӧ���������е�ע�����������˹��ࡣ�������б�żӡ���ע���ķ�ʽ��


\section{2�����⣨���룺\textbackslash section\{2������\}��}

[����������ѡ��ö��䣬�����滻֮������ΪС�ĺ����塣] 

\subsection {3�����⣨���룺\textbackslash subsection\{3������\}�� }

[����������ѡ��ö��䣬�����滻֮������ΪС�ĺ����塣]

���������֮�������

�����еı�Ű����±�Ÿ�ʽ˳�����У���(1) (2) (3), ��1) 2) 3)����٢ڢۣ���ź������ֱ����š�Ӣ����ĸ�� PPT�е���Ŀ���š�

\chapter{1������ }
[����������ѡ��ö��䣬�����滻֮������ΪС�ĺ����塣] 

\section{2������ }
\label{sec:03:01}
[����������ѡ��ö��䣬�����滻֮������ΪС�ĺ����塣] 

\subsection{3������ }
[����������ѡ��ö��䣬�����滻֮������ΪС�ĺ����塣]  
% !TeX root = ../thuthesis-example.tex

\chapter{引用文献的标注}

模板支持 BibTeX 和 BibLaTeX 两种方式处理参考文献。
下文主要介绍 BibTeX 配合 \pkg{natbib} 宏包的主要使用方法。


\section{顺序编码制}

在顺序编码制下，默认的 \cs{cite} 命令同 \cs{citep} 一样，序号置于方括号中，
引文页码会放在括号外。
统一处引用的连续序号会自动用短横线连接。

\thusetup{
  cite-style = super,
}
\noindent
\begin{tabular}{l@{\quad$\Rightarrow$\quad}l}
  \verb|\cite{zhangkun1994}|               & \cite{zhangkun1994}               \\
  \verb|\citet{zhangkun1994}|              & \citet{zhangkun1994}              \\
  \verb|\citep{zhangkun1994}|              & \citep{zhangkun1994}              \\
  \verb|\cite[42]{zhangkun1994}|           & \cite[42]{zhangkun1994}           \\
  \verb|\cite{zhangkun1994,zhukezhen1973}| & \cite{zhangkun1994,zhukezhen1973} \\
\end{tabular}


也可以取消上标格式，将数字序号作为文字的一部分。
建议全文统一使用相同的格式。

\thusetup{
  cite-style = inline,
}
\noindent
\begin{tabular}{l@{\quad$\Rightarrow$\quad}l}
  \verb|\cite{zhangkun1994}|               & \cite{zhangkun1994}               \\
  \verb|\citet{zhangkun1994}|              & \citet{zhangkun1994}              \\
  \verb|\citep{zhangkun1994}|              & \citep{zhangkun1994}              \\
  \verb|\cite[42]{zhangkun1994}|           & \cite[42]{zhangkun1994}           \\
  \verb|\cite{zhangkun1994,zhukezhen1973}| & \cite{zhangkun1994,zhukezhen1973} \\
\end{tabular}



\section{著者-出版年制}

著者-出版年制下的 \cs{cite} 跟 \cs{citet} 一样。

\thusetup{
  cite-style = author-year,
}
\noindent
\begin{tabular}{@{}l@{$\Rightarrow$}l@{}}
  \verb|\cite{zhangkun1994}|                & \cite{zhangkun1994}                \\
  \verb|\citet{zhangkun1994}|               & \citet{zhangkun1994}               \\
  \verb|\citep{zhangkun1994}|               & \citep{zhangkun1994}               \\
  \verb|\cite[42]{zhangkun1994}|            & \cite[42]{zhangkun1994}            \\
  \verb|\citep{zhangkun1994,zhukezhen1973}| & \citep{zhangkun1994,zhukezhen1973} \\
\end{tabular}

\vskip 2ex
\thusetup{
  cite-style = super,
}
注意，引文参考文献的每条都要在正文中标注
\cite{zhangkun1994,zhukezhen1973,dupont1974bone,zhengkaiqing1987,%
  jiangxizhou1980,jianduju1994,merkt1995rotational,mellinger1996laser,%
  bixon1996dynamics,mahui1995,carlson1981two,taylor1983scanning,%
  taylor1981study,shimizu1983laser,atkinson1982experimental,%
  kusch1975perturbations,guangxi1993,huosini1989guwu,wangfuzhi1865songlun,%
  zhaoyaodong1998xinshidai,biaozhunhua2002tushu,chubanzhuanye2004,%
  who1970factors,peebles2001probability,baishunong1998zhiwu,%
  weinstein1974pathogenic,hanjiren1985lun,dizhi1936dizhi,%
  tushuguan1957tushuguanxue,aaas1883science,fugang2000fengsha,%
  xiaoyu2001chubanye,oclc2000about,scitor2000project%
}。



% 其他部分
\backmatter

% 参考文献
\bibliography{ref/refs}  % 参考文献使用 BibTeX 编译
% \printbibliography       % 参考文献使用 BibLaTeX 编译

% 附录
% 本科生需要将附录放到声明之后，个人简历之前
\appendix
% % !TeX root = ../thuthesis-example.tex

\begin{survey}
\label{cha:survey}

\title{Title of the Survey}
\maketitle


\tableofcontents


本科生的外文资料调研阅读报告。


\section{Figures and Tables}

\subsection{Figures}

An example figure in appendix (Figure~\ref{fig:appendix-survey-figure}).

\begin{figure}
  \centering
  \includegraphics[width=0.6\linewidth]{example-image-a.pdf}
  \caption{Example figure in appendix}
  \label{fig:appendix-survey-figure}
\end{figure}


\subsection{Tables}

An example table in appendix (Table~\ref{tab:appendix-survey-table}).

\begin{table}
  \centering
  \caption{Example table in appendix}
  \begin{tabular}{ll}
    \toprule
    File name       & Description                                         \\
    \midrule
    thuthesis.dtx   & The source file including documentation and comments \\
    thuthesis.cls   & The template file                                   \\
    thuthesis-*.bst & BibTeX styles                                       \\
    thuthesis-*.bbx & BibLaTeX styles for bibliographies                  \\
    thuthesis-*.cbx & BibLaTeX styles for citations                       \\
    \bottomrule
  \end{tabular}
  \label{tab:appendix-survey-table}
\end{table}


\section{Equations}

An example equation in appendix (Equation~\eqref{eq:appendix-survey-equation}).
\begin{equation}
  \frac{1}{2 \uppi \symup{i}} \int_\gamma f = \sum_{k=1}^m n(\gamma; a_k) \mathscr{R}(f; a_k)
  \label{eq:appendix-survey-equation}
\end{equation}


\section{Citations}

Example\cite{dupont1974bone} citations\cite{merkt1995rotational} in appendix
\cite{dupont1974bone,merkt1995rotational}.


% 默认使用正文的参考文献样式；
% 如果使用 BibTeX，可以切换为其他兼容 natbib 的 BibTeX 样式。
\bibliographystyle{unsrtnat}
% \bibliographystyle{IEEEtranN}

% 默认使用正文的参考文献 .bib 数据库；
% 如果使用 BibTeX，可以改为指定数据库，如 \bibliography{ref/refs}。
\printbibliography

\end{survey}
       % 本科生：外文资料的调研阅读报告
% % !TeX root = ../thuthesis-example.tex

\begin{translation}
\label{cha:translation}

\title{书面翻译题目}
\maketitle

\tableofcontents


本科生的外文资料书面翻译。


\section{图表示例}

\subsection{图}

附录中的图片示例（图~\ref{fig:appendix-translation-figure}）。

\begin{figure}
  \centering
  \includegraphics[width=0.6\linewidth]{example-image-a.pdf}
  \caption{附录中的图片示例}
  \label{fig:appendix-translation-figure}
\end{figure}


\subsection{表格}

附录中的表格示例（表~\ref{tab:appendix-translation-table}）。

\begin{table}
  \centering
  \caption{附录中的表格示例}
  \begin{tabular}{ll}
    \toprule
    文件名          & 描述                         \\
    \midrule
    thuthesis.dtx   & 模板的源文件，包括文档和注释 \\
    thuthesis.cls   & 模板文件                     \\
    thuthesis-*.bst & BibTeX 参考文献表样式文件    \\
    thuthesis-*.bbx & BibLaTeX 参考文献表样式文件  \\
    thuthesis-*.cbx & BibLaTeX 引用样式文件        \\
    \bottomrule
  \end{tabular}
  \label{tab:appendix-translation-table}
\end{table}


\section{数学公式}

附录中的数学公式示例（公式\eqref{eq:appendix-translation-equation}）。
\begin{equation}
  \frac{1}{2 \uppi \symup{i}} \int_\gamma f = \sum_{k=1}^m n(\gamma; a_k) \mathscr{R}(f; a_k)
  \label{eq:appendix-translation-equation}
\end{equation}


\section{文献引用}

附录\cite{dupont1974bone}中的参考文献引用\cite{merkt1995rotational}示例
\cite{dupont1974bone,merkt1995rotational}。


\appendix

\section{附录}

附录的内容。


% 书面翻译的参考文献
% 默认使用正文的参考文献样式；
% 如果使用 BibTeX，可以切换为其他兼容 natbib 的 BibTeX 样式。
\bibliographystyle{unsrtnat}
% \bibliographystyle{IEEEtranN}

% 默认使用正文的参考文献 .bib 数据库；
% 如果使用 BibTeX，可以改为指定数据库，如 \bibliography{ref/refs}。
\printbibliography

% 书面翻译对应的原文索引
\begin{translation-index}
  \nocite{mellinger1996laser}
  \nocite{bixon1996dynamics}
  \nocite{carlson1981two}
  \bibliographystyle{unsrtnat}
  \printbibliography
\end{translation-index}

\end{translation}
  % 本科生：外文资料的书面翻译
% !TeX root = ../thuthesis-example.tex

\chapter{补充内容}

附录是与论文内容密切相关、但编入正文又影响整篇论文编排的条理和逻辑性的资料，例如某些重要的数据表格、计算程序、统计表等，是论文主体的补充内容，可根据需要设置。

附录中的图、表、数学表达式、参考文献等另行编序号，与正文分开，一律用阿拉伯数字编码，
但在数码前冠以附录的序号，例如“图~\ref{fig:appendix-figure}”，
“表~\ref{tab:appendix-table}”，“式\eqref{eq:appendix-equation}”等。


\section{插图}

% 附录中的插图示例（图~\ref{fig:appendix-figure}）。

\begin{figure}
  \centering
  \includegraphics[width=0.6\linewidth]{example-image-a.pdf}
  \caption{附录中的图片示例}
  \label{fig:appendix-figure}
\end{figure}


\section{表格}

% 附录中的表格示例（表~\ref{tab:appendix-table}）。

\begin{table}
  \centering
  \caption{附录中的表格示例}
  \begin{tabular}{ll}
    \toprule
    文件名          & 描述                         \\
    \midrule
    thuthesis.dtx   & 模板的源文件，包括文档和注释 \\
    thuthesis.cls   & 模板文件                     \\
    thuthesis-*.bst & BibTeX 参考文献表样式文件    \\
    thuthesis-*.bbx & BibLaTeX 参考文献表样式文件  \\
    thuthesis-*.cbx & BibLaTeX 引用样式文件        \\
    \bottomrule
  \end{tabular}
  \label{tab:appendix-table}
\end{table}


\section{数学表达式}

% 附录中的数学表达式示例（式\eqref{eq:appendix-equation}）。
\begin{equation}
  \frac{1}{2 \uppi \symup{i}} \int_\gamma f = \sum_{k=1}^m n(\gamma; a_k) \mathscr{R}(f; a_k)
  \label{eq:appendix-equation}
\end{equation}


\section{文献引用}

附录\cite{dupont1974bone}中的参考文献引用\cite{zhengkaiqing1987}示例
\cite{dupont1974bone,zhengkaiqing1987}。

\printbibliography


% 致谢
% !TeX root = ../thuthesis-example.tex

\begin{acknowledgements}
  衷心感谢导师×××教授和物理系××副教授对本人的精心指导。他们的言传身教将使我终生受益。

  在美国麻省理工学院化学系进行九个月的合作研究期间，承蒙 Robert Field 教授热心指导与帮助，不胜感激。

  感谢×××××实验室主任×××教授，以及实验室全体老师和同窗们学的热情帮助和支持！

  本课题承蒙国家自然科学基金资助，特此致谢。
\end{acknowledgements}


% 声明
% 本科生开题报告不需要
\statement
% 将签字扫描后的声明文件 scan-statement.pdf 替换原始页面
% \statement[file=scan-statement.pdf]
% 本科生编译生成的声明页默认不加页脚，插入扫描版时再补上；
% 研究生编译生成时有页眉页脚，插入扫描版时不再重复。
% 也可以手动控制是否加页眉页脚
% \statement[page-style=empty]
% \statement[file=scan-statement.pdf, page-style=plain]

% 个人简历、在学期间完成的相关学术成果
% 本科生可以附个人简历，也可以不附个人简历
% !TeX root = ../thuthesis-example.tex

\begin{resume}

  \section*{个人简历}

  197× 年 ×× 月 ×× 日出生于四川××县。

  1992 年 9 月考入××大学化学系××化学专业，1996 年 7 月本科毕业并获得理学学士学位。

  1996 年 9 月免试进入清华大学化学系攻读××化学博士至今。


  \section*{在学期间完成的相关学术成果}

  \subsection*{学术论文}

  \begin{achievements}
    \item Yang Y, Ren T L, Zhang L T, et al. Miniature microphone with silicon-based ferroelectric thin films[J]. Integrated Ferroelectrics, 2003, 52:229-235.
    \item 杨轶, 张宁欣, 任天令, 等. 硅基铁电微声学器件中薄膜残余应力的研究[J]. 中国机械工程, 2005, 16(14):1289-1291.
    \item 杨轶, 张宁欣, 任天令, 等. 集成铁电器件中的关键工艺研究[J]. 仪器仪表学报, 2003, 24(S4):192-193.
    \item Yang Y, Ren T L, Zhu Y P, et al. PMUTs for handwriting recognition. In press[J]. (已被Integrated Ferroelectrics录用)
  \end{achievements}


  \subsection*{专利}

  \begin{achievements}
    \item 任天令, 杨轶, 朱一平, 等. 硅基铁电微声学传感器畴极化区域控制和电极连接的方法: 中国, CN1602118A[P]. 2005-03-30.
    \item Ren T L, Yang Y, Zhu Y P, et al. Piezoelectric micro acoustic sensor based on ferroelectric materials: USA, No.11/215, 102[P]. (美国发明专利申请号.)
  \end{achievements}

\end{resume}


% 指导教师/指导小组评语
% 本科生不需要
% !TeX root = ../thuthesis-example.tex

\begin{comments}
% \begin{comments}[name = {指导小组评语}]
% \begin{comments}[name = {Comments from Thesis Supervisor}]
% \begin{comments}[name = {Comments from Thesis Supervision Committee}]

  论文提出了……

\end{comments}


% 答辩委员会决议书
% 本科生不需要
% !TeX root = ../thuthesis-example.tex

\begin{resolution}

  论文提出了……

  论文取得的主要创新性成果包括：

  1. ……

  2. ……

  3. ……

  论文工作表明作者在×××××具有×××××知识，具有××××能力，论文××××，答辩××××。

  答辩委员会表决，（×票/一致）同意通过论文答辩，并建议授予×××（姓名）×××（门类）学博士/硕士学位。

\end{resolution}


% 本科生的综合论文训练记录表（扫描版）
% \record{file=scan-record.pdf}

\end{document}
